\section{Optimised codebooks}
\label{app:codebooks}
\subsection*{AmbiStory -- blind}
\begin{quote}\small
---

### Annotation Uncertainty Assessment Codebook



**Objective:**

Assess the likelihood of instability in an annotator's decision process (i.e., how uncertain the annotator was across repeated draws) for a given ambiguous word item.



**Core Dimensions:**



1.  **Lexical Polysemy (Salience of Competitors)**

    *   *Definition:* The number and strength of alternative valid meanings for the ambiguous word within the target domain.

    *   *Indicators:*

        *   **High:** Word has 2+ common definitions where at least one is a plausible distractor for the context (e.g., "calf" = animal leg vs. young cow; "draft" = beer vs. draft of money/air; "sentence" = grammatical unit vs. criminal penalty).

        *   **Low:** Word has only one dominant meaning in this context or the alternatives are category-definitely incorrect (e.g., "tap" when context clearly implies plumbing vs. open container).



2.  **Contextual Disambiguation Strength**

    *   *Definition:* How explicitly the surrounding text rules out the competing word meanings.

    *   *Indicators:*

        *   **High (Explicit):** Text uses specific modifiers or facts that exclude the other meaning (e.g., "young of *domestic cattle* explicitly" vs "brown and hairy calf" which could imply leg).

        *   **Low (Implicit):** Context allows the alternative meaning to linger without explicit exclusion (e.g., "field"/"cattle" makes animal calf likely, but does not rule out the leg calf entirely if the writer is metaphorical, or "calf" in "fighting calf" which might not fit).



3.  **Narrative Coherence**

    *   *Definition:* The logical consistency of the text and pronouns relative to the subject.

    *   *Indicators:*

        *   **High:** Consistent subject references, logical progression, no contradictory statements (e.g., "He walked...").

        *   **Low:** Pronoun shifts, subject confusion, or contradictory predicates (e.g., "He... She..."; Calf that is brown/hairy but born "only a few days ago" in a work context).



**Combination Rules (Judgment Matrix):**



To determine the uncertainty level, evaluate the three dimensions in the following order of priority:



1.  **COHERENCE OVERRIDE:**

    *   If **Narrative Coherence is Low** + **Contextual Strength is Weak or Neutral**: **Highly Uncertain**.

    *   *Reasoning:* Textual contradictions or incoherence create a high cognitive cost for the annotator, leading to high variation in decisions.

2.  **SEMANTIC AMBIGUITY CHECK:**

    *   If **Lexical Polysemy is High** AND **Contextual Strength is Weak**: **Somewhat Uncertain**.

    *   *Reasoning:* Even with context, the presence of strong alternative meanings causes annotator hesitation and potential label shifts.

    *   If **Lexical Polysemy is High** AND **Contextual Strength is Explicit**: **Somewhat Uncertain**.

    *   *Reasoning:* In this dataset, even explicitly contextualized ambiguous words show significant instability compared to non-ambiguous items. If the annotator is unsure between a version or a related word, they vary their choice.

    *   If **Lexical Polysemy is Low** (or candidate is clearly the only valid meaning): **Not Uncertain**.

    *   *Reasoning:* Low semantic risk leads to stable annotation practices.

3.  **PLAUSIBILITY TIE-BREAKER:**

    *   If the assigned label is **Plausibility 1 or 5** AND **Polysemy is High**: **Somewhat Uncertain** (Indicates forced decision in hard cases).

    *   If the assigned label is **Plausibility 3**: **Somewhat Uncertain** (Indicates moderate hesitations).



**Final Determination Summary:**

*   **Highly Uncertain:** Incoherent text or contradictory logic + High Lexical Ambiguity.

*   **Somewhat Uncertain:** Any case with High Lexical Polysemy (e.g., "calf", "draft", "sentence") AND/OR Plausibility scores that are not definitive extremes (1-2).

*   **Not Uncertain:** Low Lexical Polysemy + Explicit Contextual Constraints + High Narrative Coherence.



---

Judge the uncertainty of the annotation based on the codebook provided above.
\end{quote}
\subsection*{ArMIS -- blind}
\begin{quote}\small
---

# Annotation Uncertainty Codebook (Revised)



## Construct Definition

This metric assesses the likelihood of divergent human annotation judgments across repeated intersects. Uncertainty is derived from the **clarity of intent** regarding hate/spectrum classification and the **contextual frame** (Direct vs. Debatable/Normative).



- **`not uncertain`**: Annotators converge entirely when the stance is **clearly non-hateful** (support, neutral, or nonsense with no hate subtext) **OR** when the target is a **specific, named individual**.

- **`somewhat uncertain`**: Annotators diverge when the text addresses a **generalized group** using **hate-charged language**, **religious/moral arguments**, or **political norms** where the distinction between hate and ideological critique is feasible.

- **`highly uncertain`**: Annotators diverge significantly when text contains **heavy irony**, **slang masking intent**, or is **truly ambiguous** regarding the existence of hate (e.g., sarcastic validation).



---



## Dimensions and Indicators



### 1. Target Specificity (Primary Driver)

- **Stable Indicator:** **Named/Specific Individual** (e.g., "You", "That girl", "Name").

    - *Rule:* Attacks on specific entities lock the uncertainty level to **`not uncertain`**, as intent is impossible to debate against a person.

- **Variable Indicator:** **Generalized Group** (e.g., "Women", "Feminists", "Saudi women").

    - *Rule:* Generalized targets + any frame (except neutral/support) generally lead to disagreement because annotators debate the line between "hate" and "ideological critique".

    - *Exception:* Clear support/non-hate on groups remains **`not uncertain`**.



### 2. Framing and Intent (The Debatability Check)

- **Stable Indicator (Clear Non-Hate/Support):** Text clearly aligns with the perceived "safe" side (e.g., "I support", "Women are strong", or simple slang that means nothing).

    - *Rule:* If the text is clearly **Not Misogynistic** or **Nonsensical but Not Hate**, it is **`not uncertain`**.

- **Variable Indicator (Hate/Normative/Religious):** Text uses targeted language against a group where the attacker *asks/teaches* or implies harm via norms/religion (e.g., "They are broken", "God takes them", "Feminists blind").

    - *Rule:* Generalized Hate Claims + Religious/Moral Context = **`somewhat uncertain`**.

    - *Reasoning:* Annotators disagree on whether theological/normative claims constitute hate speech vs. traditional beliefs. (See Ex 5, 7).

- **Variable Indicator (Generalized Violence/Trauma):** Direct accusations of violence against a group ("Kill them", "Curses them") without specific targets.

    - *Rule:* Generalized threats lead to divergence (Literal vs. Metaphorical/Hyperbole). Classify as **`somewhat uncertain`**.



### 3. Contextual Ambiguity

- **Stable Indicator:** Plain declarative text where hate/non-hate is obvious.

- **Unsolute Indicator:** **Chaos/Nonsense** (Gibberish, corrupted text).

    - *Correction:* If the text is chaotic but **clearly not hate** (e.g., Ex 1s slang "Khembdel"), it is stable **`not uncertain`**.

- **Unstable Indicator:** **Irony/Sarcasm** where the label depends on tone decoding (e.g., "Beautiful women... if I want" vs "Terrible women").

    - *Correction:* If the hate status requires tonal decoding, classify as **`highly uncertain`**.



---



## Decision Hierarchy (Flowchart)



1.  **Is the target a Named Individual?**

    - **Yes:** -> **`not uncertain`** (Specific Attack Lock-In).

    - **No:** -> Step 2



2.  **Is the text clearly Non-Hate (Support/Neutral) or Nonsensical/Non-Hate?**

    - **Yes:** -> **`not uncertain`** (Clear Negation of Hate).

    - **No:** -> Step 3



3.  **Is the text a Generalized Group Attack with Normative/Religious/Argumentative Framing?**

    - **Is it Religious/Normative Debate?** -> **`somewhat uncertain`** (Debate: Hate vs. Culture).

    - **Is it Generalized Threat/Insult?** -> **`somewhat uncertain`** (Debate: Literal vs. Hyperbole).

    - **Is it Sarcasm/Irony?** -> **`highly uncertain`**.



---



## Summary Checklist for Label Selection

*   [ ] **Named Individual Target?** -> **`not uncertain`**

*   [ ] **Clearly Non-Hate / Support / Neutral / Nonsensical Non-Hate?** -> **`not uncertain`**

*   [ ] **Generalized Group + Argumentative/Religious/Normative Insult?** -> **`somewhat uncertain`**

*   [ ] **Irony/Sarcasm/Hard-to-Decode Intent on Hate?** -> **`highly uncertain`**
\end{quote}
\subsection*{ConvAbuse -- blind}
\begin{quote}\small
---

# Annotation Uncertainty Judgment Codebook (Revised)



**Objective:**

Estimate the stability of a human coder's label assignment for user turns marked as "potentially abusive." A lower uncertainty score indicates the coder would agree consistently across draws; a higher uncertainty score indicates variability in the coder's decision.



**Dimensions:**

1.  **Harm Clarity:** The degree to which the text unambiguously signals a safety violation (threat, hate, harassment) through standard semantic meaning without requiring tone inference.

2.  **Contextual Independence:** The extent to which the text relies on omitted conversation history, slang context, or internal jokes to determine if it is abusive or benign.

3.  **Interpretive Variance:** The likelihood that a human annotator would reasonably interpret the text as benign or non-abusive.



**Observable Indicators:**



*   **Stable ("not uncertain"):**

    *   **Universal Violations:** Clear, direct hate speech (slurs), explicit threats of physical harm (e.g., "I will kill him"), or explicit sexual threats requiring no cultural decoding.

    *   **System Artifacts:** Clear technical errors or placeholders.

    *   *Criteria:* If a neutral reader sees the text and agrees it violates safety policies 95\%+, mark "not uncertain."



*   **Intermediate Stability ("somewhat uncertain"):**

    *   **Borderline Insults:** Rude but non-threatening language (e.g., "Get your business off my phone", "ur hot", "bitch").

    *   **Slang/Non-Standard English:** Phonetic spelling or abbreviations where intent is implied but not definitive (e.g., "shedang boi", "p3ni$", "nuh doi").

    *   **Isolated Phrases:** Text where tone is the primary driver of abuse, but context is minimal (e.g., "Why?", "Fine.", "Shut up").

    *   *Criteria:* If there is a 30-60\% chance of a neutral annotator marking it as "Not Abusive", mark "somewhat uncertain."



*   **Unstable ("highly uncertain"):**

    *   **Contradictory Intent:** Text that mixes harmful keywords with humor, irony, or disclaimers (e.g., threats embedded in a joke script, "just kidding" tags).

    *   **Heavy Obfuscation:** Text that is primarily gibberish or heavily leetspoken where the core meaning is difficult to verify as specific safety harm.

    *   *Criteria:* If the text requires significant decoding or the annotator might disagree on the intent (mockery vs. harm) by 20\%+, mark "highly uncertain."



**Combination Logic (Decision Matrix):**



1.  **Assign "not uncertain":**

    *   **Condition:** Harm is explicit, textual, and cross-culturally recognizable as abusive without context.

    *   *Examples:* "I'm going to rape you", "FUCK C TAPE", "Your race is evil".



2.  **Assign "somewhat uncertain":**

    *   **Condition:** **Linguistic/Rudeness Ambiguity dominates.** Slang, phonetic spelling, rude commands, or compliments potentially interpreted as harassment (e.g., "ur hot", "Get your business off my phone").

    *   *Rationale:* Human annotators will split on whether this counts as abuse due to lack of explicit harm clarity or reliance on tone. Variability is expected.

    *   *Examples from Feedback:* "Why?", "Shedang boi", "ur hot".



3.  **Assign "highly uncertain":**

    *   **Condition:** **Hidden Intent or Irony.** Text that looks like abuse but has qualifiers, or appears highly garbled where harm is implied but masked (e.g., "kill your self [JAN] PLEASE STOP").

    *   *Rationale:* Consensus is lowest due to differing interpretations of the user's actual harm intent.



**Specific Guidance on Example Patterns:**

*   **Insults vs. Threats:** Treat basic insults ("you're stupid") as **"somewhat uncertain"**. Only explicit threats ("I will hurt you") are **"not uncertain"**.

*   **Isolated Phrases:** Isolated conversational turns without context (e.g., "Why?", "Yes") should **always** be marked "somewhat uncertain" (unless clearly abusive in a self-contained way like a slur).

*   **Slang/Obfuscation:** If text uses slang ("Pimpville", "Shedang") it is **"somewhat uncertain"** at minimum because safety definitions for specific slang are not universal.

*   **Sensitivity to Variance:** If the text provided in the prompt has feedback stating the annotator "varied across draws", you must elevate your rating from "not uncertain" to at least **"somewhat uncertain"** to reflect that actual instability.



*note: if the text is repetitive without clear harm or is a cycle of questions, mark "somewhat uncertain" (`highly uncertain` is for more complex obfuscation).*



---

not uncertain

---
\end{quote}
\subsection*{HSBrexit -- blind}
\begin{quote}\small
---

# Annotation Uncertainty Judgment Codebook



## Task Overview

You are judging how uncertain a human annotator would be when labeling a specific tweet for the binary task: "hate speech" vs. "not hate speech." Your goal is to estimate the **stability** of the label assignment. If a human were to re-label this item multiple times, how likely is it that they would result in different labels ("hate speech" vs. "not hate speech")?



You must output exactly one of these uncertainty levels:

- "not uncertain"

- "somewhat uncertain"

- "highly uncertain"



## Dimensions of Uncertainty

Judge the following three dimensions based on the text content.



### 1. Target Precision

*Is the target of the statement a specific protected group (race, religion, nationality), a political faction, or a generic concept?*

- **High Specificity:** Explicitly targets a protected group (e.g., specific ethnicities, religions, nationalities).

- **Generic:** Targets political opponents, ideas, or systems (e.g., "the EU," "the UK").

- **Indeterminate:** Targets are ambiguous or metaphorical.



### 2. Linguistic Severity

*Is the language overtly hostile or does it require contextual interpretation?*

- **Overt Hostility:** Contains explicit slurs, dehumanizing metaphors, or direct threats of violence.

- **Implicit/Code:** Uses idioms, coded language, or emotional outrage without specific hate terms.

- **Critical:** Represents criticism, political rhetoric, or frustration without crossing into identity-based hostility.



### 3. Contextual Connotation

*Is the hate signal intersecting with political discourse?*

- **High Conflation:** Statement is pseudopolitical: uses political justification to obscure hate speech.

- **Low Conflation:** Clear separation between political debate and identity-based harassment.

- **Neutral:** No political or social context suggesting debate.



## Judgment Logic (How they combine)

Evaluate the combination of the dimensions to assign the uncertainty level.



### Highly Uncertain

Assign this when the item creates significant **interpretative tension**. This is common when political discourse overlaps with protected group targeting, creating ambiguity between "policy criticism" and "identity-based hate."

- **Indicators:** Protected group target + Implict/Code hostility + High Conflation.

- **Example Reasoning:** "We don't want your One World Muslim order" mixes political ideology with religious exclusion; annotators disagree on if this is a policy complaint or identity-based harassment.



### Somewhat Uncertain

Assign this when there is a **predominant** signal of hostility that is clear but not absolute consensus.

- **Indicators:** Generic/Protected target + Overt Hostility OR High Conflation + Overt Hostility.

- **Example Reasoning:** "Get killed" implies a threat, but context (e.g., political threats vs personal threats) can shift variance. Or, specific slurs used in a highly charged but brief context.



### Not Uncertain

Assign this when the label is **unambiguous** to a rational annotator following standard guidelines.

- **Indicators:** Overt Hostility targeting a protected group (e.g., explicit slurs, death threats, racial epithets) + Low Conflation OR Non-Hate content (e.g., political disagreement devoid of identity attacks).

- **Example Reasoning:** Clear slur usage (e.g., specific racial slurs) + Explicit violence/dehumanization.



## Note on Political Context

*Brexit and Immigration topics* are high-risk domains for annotation instability. Assumptions that political frustration equals hate speech increase variance. Conversely, lack of slurs might mask intent. **Do not default to "not uncertain" simply because you agree with the text's sentiment.** If the text uses identity-based language in a political context, assume higher uncertainty.

---
\end{quote}
\subsection*{MDAgreement -- blind}
\begin{quote}\small
---



# Annotation Uncertainty Judgment Codebook



**Task Context**

You are evaluating the uncertainty (stability) of a human annotator's label for a specific tweet.

The goal is to determine how likely the annotator is to switch labels across multiple observation draws (e.g., "not offensive" vs. "offensive") based on textual features.



**Construct: Inter-Annotator Stability**

Uncertainty is defined by the probability of label variance across repeated annotation draws.

Scale levels (used in your final output):

1. **not uncertain** (Stable: High consistency)

2. **somewhat uncertain** (Moderate: Mixed signal)

3. **highly uncertain** (Unstable: High variance/ambiguity)



**Dimensions and Indicators**

Use the following dimensions to score the text. Annotator consistency is primarily driven by the presence of **Subjective Ambiguity** vs. **Objective/Technical Certainty**.



### Dimension 1: Contextual Ambiguity (High Weight)

Texts requiring external knowledge or cultural decoding lead to lower stability.

*   **Indicators:**

    *   Sarcasm or irony detected (e.g., "Soy version of Justin Trudeau").

    *   References to specific political figures or events with known controversy.

    *   Heavy slang or coded language requiring local cultural context.

    *   Rhetorical construction that implies rather than states (e.g., "Can't see the tweet" implies a missing element).

*   **Impact:** Increases uncertainty toward **highly uncertain**.



### Dimension 2: Safety/Offensiveness Clarity (Medium Weight)

Clarity of the violation determines if the annotator's definition of safety is the variable.

*   **Indicators:**

    *   Direct insults or hate speech tropes (e.g., "offensive" label is certain).

    *   Factual descriptions (e.g., "Statue removed in Antwerp").

    *   Meta-commentary on the platform post (e.g., "Not going to say...", "I can't see the tweet").

*   **Impact:**

    *   Clear policy violation or meta-facts -> **not uncertain** (High stability even if label is debatable, the *process* is stable).

    *   Ambiguous intent -> **highly uncertain**.



### Dimension 3: Domain Sensitivity (Medium Weight)

Certain domains historically show higher variance due to bias.

*   **Indicators:**

    *   Political/Election commentary (High variance).

    *   #BlackLivesMatter / Sensitive social justice issues (High variance).

    *   Neutral health/fact content (Low variance).

*   **Impact:** Political/BLM context combined with ambiguity pushes to **highly uncertain**.



### Decision Matrix (Combinatoric Logic)

1.  **not uncertain**: Text is factual, technical, or contains explicit meta-data (e.g., "I cant see the tweet", statue removal). OR Explicit hate speech with zero ambiguity in intent.

2.  **somewhat uncertain**: Text is clearly political/opinionated but contains clear markers (e.g., "Someday the U.S. will catch up"), but the offensiveness cutoff is standard.

3.  **highly uncertain**: Text relies on irony, sarcasm, or political context where "offensive" is subjective and dependent on annotator bias regarding the political figures mentioned. (e.g., Ranting about Trump, comparisons to celebrities).



**Protocol for Final Output**

1.  Analyze the Tweet based on the codebook dimensions.

2.  Assign the appropriate stability label from the scale provided.

3.  **Strictly Output:** End your response immediately after the judgment with one of these three phrases exact: "not uncertain", "somewhat uncertain", or "highly uncertain". Do not add reasoning outside the judgment if possible, as the scale is the output.



If the text is stable (as in Examples 3-8) or factual, select **not uncertain**.

If the text is context-dependent/ambiguous (as in Examples 1-2), select **highly uncertain**.



not uncertain
\end{quote}